
\documentclass[journal,11pt]{article}

\usepackage[utf8]{inputenc}
\usepackage{amsmath,amssymb,amsfonts,amsthm,bm}
\usepackage{graphicx}
\usepackage{booktabs}
\usepackage{multirow}
\usepackage{algorithm}
\usepackage{algpseudocode}
\usepackage{hyperref}
\usepackage{geometry}
\geometry{margin=25mm}

\title{Support Tensor Theory for Additive Manufacturing Orientation Optimization}
\author{InHwan Sul et al.}
\date{}

\newtheorem{definition}{Definition}
\newtheorem{theorem}{Theorem}
\newtheorem{proposition}{Proposition}

\begin{document}

\maketitle

\begin{abstract}
This paper proposes a novel tensor-based framework for rapid build orientation estimation in additive manufacturing.
The proposed Support Tensor converts the nonlinear support-generation problem into an eigenvalue problem on a symmetric positive semidefinite matrix.
The method enables rapid identification of candidate build directions directly from STL geometry without exhaustive directional search.
\end{abstract}

\textbf{Keywords:}
Additive Manufacturing, Build Orientation, Support Structure, Eigenvalue Analysis, PCA, Tensor Methods

\section{Introduction}

Build orientation strongly influences support volume, build time, material consumption, surface quality, and post-processing cost in additive manufacturing.

Current approaches typically rely on:

\begin{enumerate}
\item exhaustive directional sampling,
\item genetic algorithms,
\item topology-aware search,
\item commercial slicer simulation.
\end{enumerate}

These approaches are computationally expensive.

This study proposes a new mathematical descriptor called the \emph{Support Tensor}.
The key idea is to represent support tendency through a second-order tensor derived from surface-normal distributions.

\section{Related Work}

PCA has long been used for geometric analysis.

For a point cloud,

\begin{equation}
C=\sum_i (\mathbf{x}_i-\bar{\mathbf{x}})(\mathbf{x}_i-\bar{\mathbf{x}})^T
\end{equation}

and eigenvectors of $C$ define principal directions.

However, support generation depends primarily on surface orientation rather than point distribution.

\section{Support Functional}

Let

\begin{itemize}
\item $A_i$ : triangle area
\item $\mathbf m_i$ : unit outward normal
\item $\mathbf n$ : build direction
\end{itemize}

The support tendency is approximated by

\begin{equation}
S(\mathbf n)
=
\sum_i
A_i
H(-\mathbf m_i\cdot \mathbf n)
(-\mathbf m_i\cdot \mathbf n).
\end{equation}

This function is nonlinear and difficult to optimize directly.

\section{Quadratic Approximation}

Consider

\begin{equation}
\widetilde S(\mathbf n)
=
\sum_i
A_i
(\mathbf m_i\cdot\mathbf n)^2.
\end{equation}

Since

\begin{equation}
(\mathbf m_i\cdot\mathbf n)^2
=
\mathbf n^T
(\mathbf m_i\mathbf m_i^T)
\mathbf n,
\end{equation}

we obtain

\begin{equation}
\widetilde S(\mathbf n)
=
\mathbf n^T M \mathbf n.
\end{equation}

\begin{definition}[Support Tensor]
\begin{equation}
M
=
\sum_i
A_i
\mathbf m_i\mathbf m_i^T.
\end{equation}
\end{definition}

\section{Theoretical Properties}

\begin{proposition}
The Support Tensor is symmetric.
\end{proposition}

\textbf{Proof}

Each term $\mathbf m_i\mathbf m_i^T$ is symmetric.

Therefore

\begin{equation}
M^T=M.
\end{equation}

\hfill $\square$

\begin{proposition}
The Support Tensor is positive semidefinite.
\end{proposition}

\textbf{Proof}

For any vector $\mathbf x$,

\begin{equation}
\mathbf x^T M \mathbf x
=
\sum_i
A_i(\mathbf x^T\mathbf m_i)^2
\ge 0.
\end{equation}

\hfill $\square$

\begin{theorem}
The minimum of

\begin{equation}
\mathbf n^T M \mathbf n
\end{equation}

under

\begin{equation}
||\mathbf n||=1
\end{equation}

is the smallest eigenvalue of $M$.
\end{theorem}

\textbf{Proof}

Directly follows from the Rayleigh Quotient theorem.

\hfill $\square$

\section{Weighted Support Tensor}

A generalized model is

\begin{equation}
M_w
=
\sum_i
w_iA_i
\mathbf m_i\mathbf m_i^T.
\end{equation}

where

\begin{equation}
w_i
=
f(q_i,d_i,r_i).
\end{equation}

Possible weighting factors:

\begin{itemize}
\item cavity influence
\item overhang sensitivity
\item support accessibility
\item local manufacturing penalty
\end{itemize}

\section{Cavity-Aware Extension}

Let

\begin{equation}
q_i\in[0,1]
\end{equation}

be a cavity score.

Then

\begin{equation}
w_i=1+\alpha q_i.
\end{equation}

This gives

\begin{equation}
M_{cavity}
=
\sum_i
(1+\alpha q_i)
A_i
\mathbf m_i\mathbf m_i^T.
\end{equation}

The eigenvectors become sensitive to internal voids and trapped regions.

\section{Algorithm}

\begin{algorithm}
\caption{Support Tensor Orientation Estimation}
\begin{algorithmic}
\For{each triangle}
\State Compute area
\State Compute unit normal
\State Update tensor
\EndFor
\State Eigen decomposition
\State Select minimum-eigenvalue eigenvector
\State Validate with support simulation
\end{algorithmic}
\end{algorithm}

\section{Experimental Design}

Benchmark models:

\begin{enumerate}
\item Stanford Bunny
\item Mechanical Bracket
\item Turbine Blade
\item Thermal Manikin Components
\item Apparel-related 3D Models
\end{enumerate}

Metrics:

\begin{equation}
E_{angle}
=
\cos^{-1}
(|\mathbf n_{pred}\cdot\mathbf n_{best}|)
\end{equation}

and

\begin{equation}
E_{support}
=
\frac{
S(\mathbf n_{pred})
-
S(\mathbf n_{best})
}
{
S(\mathbf n_{best})
}.
\end{equation}

\section{Discussion}

The proposed approach has several advantages:

\begin{enumerate}
\item Extremely low computational cost.
\item Closed-form eigenvalue solution.
\item Easy STL implementation.
\item Extendable to cavity-aware optimization.
\end{enumerate}

Potential limitations include:

\begin{enumerate}
\item inability to fully model support visibility,
\item slicer-specific rules,
\item self-shadowing effects.
\end{enumerate}

\section{Future Research}

Future work includes:

\begin{itemize}
\item support visibility tensors,
\item graph-based cavity descriptors,
\item thermal distortion weighting,
\item multi-objective optimization,
\item topology-aware support tensors.
\end{itemize}

\section{Conclusion}

Support Tensor Theory transforms orientation optimization into a tractable eigenvalue problem.
The framework provides a new connection between tensor analysis, PCA, and additive manufacturing orientation selection.

\bibliographystyle{plain}

\begin{thebibliography}{99}

\bibitem{jolliffe}
Jolliffe, I. T.
Principal Component Analysis.
Springer.

\bibitem{gibson}
Gibson, I.
Additive Manufacturing Technologies.
Springer.

\bibitem{mardia}
Mardia, K.
Directional Statistics.
Wiley.

\end{thebibliography}

\end{document}
